\chapter{Gigantism}
\index{Gigantism}

\section*{Definition and Overview}
Gigantism is a rare condition of abnormal, excessive growth that occurs when a child or adolescent has too much growth hormone before their growth plates (epiphyses) have closed. It is almost always caused by a benign tumour of the pituitary gland that secretes excess growth hormone. Because the bones are still growing, affected children grow unusually tall and large, and organs also enlarge. When the same excess growth hormone occurs after growth has finished, it is called acromegaly. Gigantism is important to recognise early because the tumour can also press on surrounding brain structures and because untreated excess growth hormone leads to serious health problems. Treatment, usually surgery and sometimes medicines or radiotherapy, can control the condition.

\section*{Epidemiology}
Gigantism is extremely rare, affecting fewer than one in a million children. It is even less common than acromegaly, its adult counterpart, which affects around 50--70 people per million. Gigantism is slightly more common in boys, and most cases are diagnosed between the ages of 10 and 15, when rapid growth becomes obvious. Because the condition is so rare, diagnosis is often delayed, and the average delay between the start of excessive growth and diagnosis can be several years. In a minority of families, gigantism occurs as part of inherited conditions such as multiple endocrine neoplasia type 1 or familial isolated pituitary adenoma.

\section*{Aetiology and Pathophysiology}
Gigantism results from overproduction of growth hormone (GH) during the growing years.

\begin{itemize}
    \item \textbf{Pituitary adenoma:} a benign tumour of the pituitary gland secretes excess growth hormone in over 90\% of cases
    \item \textbf{Growth hormone excess:} high GH stimulates the liver to produce insulin-like growth factor 1 (IGF-1), which drives growth of bone and soft tissue
    \item \textbf{Open growth plates:} because growth plates are still open in children and adolescents, excess GH causes elongation of bones and extreme height
    \item \textbf{Organ enlargement:} internal organs, including the heart, also enlarge
    \item \textbf{Tumour effects:} the adenoma can compress the optic nerves, the pituitary gland, and surrounding brain tissue
    \item \textbf{Genetic syndromes:} rare inherited conditions can predispose to the pituitary tumours that cause gigantism
\end{itemize}

\section*{Clinical Features}
The hallmark is growth that is far faster and taller than expected for the child's age and family.

\begin{itemize}
    \item Height and growth velocity well above the 97th percentile for age and sex
    \item Large hands and feet, with rapid shoe size increases
    \item Coarse facial features and a prominent jaw
    \item Enlarged forehead, nose, lips, and tongue
    \item Headaches, which may indicate pressure from the tumour
    \item Visual problems, including loss of peripheral vision, from optic nerve compression
    \item Excessive sweating, fatigue, and joint aches
    \item Delayed puberty in some cases
    \item Enlargement of internal organs, sometimes causing heart problems
\end{itemize}

\section*{Differential Diagnosis}
Other causes of excessive growth in childhood should be considered before diagnosing gigantism.

\begin{itemize}
    \item Constitutional tall stature: children who are simply tall because of their parents' height
    \item Familial tall stature syndromes
    \item Precocious puberty, which causes a temporary growth spurt
    \item Marfan syndrome and other connective tissue disorders
    \item Klinefelter syndrome and other chromosomal conditions
    \item Sotos syndrome and other overgrowth syndromes
    \item Obesity-related growth acceleration
    \item Thyroid disorders and other endocrine conditions
\end{itemize}

\section*{When to Seek Medical Care}
Parents should seek medical review if a child's growth is far outpacing their peers, if shoes and clothes need replacing at an unusually rapid rate, or if there is any delay or acceleration of puberty. Early referral to a paediatric endocrinologist improves outcomes.

\textbf{Contact your local emergency services for:}
\begin{itemize}
    \item Sudden severe headache with vomiting, visual loss, or confusion, which may indicate bleeding into the pituitary tumour
    \item Sudden loss of vision
    \item Severe chest pain or breathlessness, which may reflect cardiac complications
    \item Any acute neurological symptoms such as weakness or seizures
\end{itemize}

\section*{Diagnosis and Investigations}
The diagnosis is confirmed by blood tests and imaging of the pituitary gland.

\begin{itemize}
    \item \textbf{Growth records:} plotting height and weight on growth charts documents excessive growth velocity
    \item \textbf{Blood tests:} elevated growth hormone and insulin-like growth factor 1 (IGF-1) levels are the key findings
    \item \textbf{Glucose tolerance test:} in gigantism, growth hormone does not suppress normally after a glucose drink
    \item \textbf{MRI of the pituitary:} identifies the adenoma and its size and position
    \item \textbf{Visual field testing:} checks for optic nerve compression from the tumour
    \item \textbf{Assessment for other hormone problems:} the tumour can disrupt other pituitary hormones, so they are checked
    \item \textbf{Genetic testing:} considered when there is a family history or features of an inherited syndrome
\end{itemize}

\section*{Management}
Treatment aims to remove or control the growth hormone-secreting tumour and to restore normal growth hormone levels.

\begin{itemize}
    \item \textbf{Surgery:} transsphenoidal surgery (through the nose) to remove the adenoma is the first-line treatment and can cure many cases
    \item \textbf{Medicines:} somatostatin analogues (such as octreotide) and growth hormone receptor antagonists reduce or block the effects of growth hormone
    \item \textbf{Radiotherapy:} used for tumours that cannot be fully removed or are resistant to medicines; effects develop slowly
    \item \textbf{Multidisciplinary care:} paediatric endocrinologists, neurosurgeons, and ophthalmologists work together
    \item \textbf{Monitoring:} regular blood tests, growth measurements, and imaging track the response to treatment
    \item \textbf{Management of other hormone deficiencies:} replacement of thyroid, adrenal, or sex hormones if the tumour or surgery damages the pituitary
    \item \textbf{Psychosocial support:} children and families benefit from support with the physical and emotional effects of extreme height
\end{itemize}

\section*{Complications}
\begin{itemize}
    \item Permanent vision loss from optic nerve compression
    \item Heart enlargement, high blood pressure, and other cardiovascular problems
    \item Diabetes and glucose intolerance from growth hormone excess
    \item Joint and bone problems, including arthritis and spinal curvature
    \item Obstructive sleep apnoea from enlarged airways and tissues
    \item Damage to normal pituitary function, requiring lifelong hormone replacement
    \item Effects of extreme stature on schooling, social life, and self-esteem
\end{itemize}

\section*{Prognosis and Outcomes}
With early diagnosis and treatment, the outlook for gigantism has improved greatly. Surgery cures a substantial proportion of children, and those who cannot be cured can usually be well controlled with medicines. Once growth hormone levels are controlled, growth returns to a normal trajectory, and adult height, while often still tall, is much closer to normal. Lifelong monitoring is needed because the tumour can recur and because long-term health requires normal hormone levels. Untreated gigantism is associated with significant early mortality from cardiovascular disease, which is why prompt treatment matters.

\section*{Prevention and Self-Care}
\begin{itemize}
    \item Track your child's height, weight, and growth velocity at routine health checks
    \item Seek review early if a child's growth rate is dramatically faster than their peers or siblings
    \item Attend all follow-up appointments and blood tests after treatment
    \item Take prescribed hormone treatments exactly as directed
    \item Maintain regular vision checks if the pituitary tumour was near the optic nerves
    \item Support the child's emotional wellbeing, and seek counselling if stature is affecting self-esteem
    \item Watch for symptoms such as headaches, visual changes, or sweating, and report them promptly
    \item If a family member has a pituitary condition, discuss genetic testing with a specialist
\end{itemize}

\section*{Sources and Further Reading}
The following reputable sources provide further information for healthcare students and patients seeking reliable, current advice about gigantism.

\begin{itemize}
    \item Pituitary Foundation (Australia). \textit{Information about growth hormone excess.} https://www.pituitary.org.au/
    \item Australian Rheumatology Association. \textit{Pituitary and growth disorders information.} https://rheumatology.org.au/
    \item Healthdirect Australia. \textit{Gigantism and acromegaly.} https://www.healthdirect.gov.au/
    \item NHS. \textit{Acromegaly and gigantism.} https://www.nhs.uk/
    \item National Institute of Diabetes and Digestive and Kidney Diseases. \textit{Acromegaly and gigantism.} https://www.niddk.nih.gov/
    \item MSD Manual. \textit{Gigantism and acromegaly.} https://www.msdmanuals.com/
\end{itemize}
